\section{Sliding Window Sum}
\label{sec:cses-sw-sum}

\problemstatement{Given an array of $n$ integers and a window size $k$,
output the sum of every contiguous window of size $k$ (there are $n-k+1$
of them).}

\begin{patternbox}
Sum is \textbf{invertible}, so the window answer updates in $O(1)$: when
the window slides one step, add the new right element and subtract the old
left element. This is the simplest instance of the add/remove/report
template.
\end{patternbox}

\subsection*{A running total}

Compute the sum of the first window directly. Then each slide changes the
window by exactly two elements: one enters on the right, one leaves on the
left. Because addition has an inverse, the new sum is
$\textit{sum}+a[\text{new}]-a[\text{old}]$ -- no re-summation. (Equivalently,
a prefix-sum array answers each window as $\textit{pre}[i+k]-\textit{pre}
[i]$.)

\begin{ideabox}[Invertible Aggregates Slide for Free]
Any aggregate with an inverse -- sum, XOR, product over a field -- lets the
window update in constant time by undoing the departing element and
applying the arriving one. Reach for this before any heavier structure.
\end{ideabox}

\subsection*{Algorithm}

\begin{enumerate}
  \item Sum the first $k$ elements.
  \item For each subsequent window, add $a[i]$ and subtract $a[i-k]$;
        record the sum.
\end{enumerate}

\subsection*{Dry run}

\dryrun{$a=[1,3,2,5,1]$, $k=3$.}

\begin{itemize}
  \item Window $[1,3,2]$: sum $6$.
  \item Slide: $+5-1\Rightarrow[3,2,5]$ sum $10$.
  \item Slide: $+1-3\Rightarrow[2,5,1]$ sum $8$. Output $6,10,8$.
\end{itemize}

\subsection*{C++ implementation}

\begin{lstlisting}
#include <bits/stdc++.h>
using namespace std;

int main() {
    int n, k;
    cin >> n >> k;
    vector<long long> a(n);
    for (auto& x : a) cin >> x;

    long long sum = 0;
    for (int i = 0; i < k; i++) sum += a[i];
    cout << sum << " ";
    for (int i = k; i < n; i++) {
        sum += a[i] - a[i - k];               // add new, drop old
        cout << sum << " ";
    }
    cout << "\n";
    return 0;
}
\end{lstlisting}

\begin{complexitybox}
\textbf{Time Complexity.} $O(n)$ -- one pass with $O(1)$ per window.
\textbf{Space Complexity.} $O(1)$ beyond the input.
\end{complexitybox}

\begin{takeawaybox}
When the window statistic is invertible, sliding is a constant-time
add-and-subtract. Establish this baseline first; the harder window problems
are exactly those where the statistic (min, mode, median) has no cheap
inverse and demands a supporting structure.
\end{takeawaybox}
